\documentclass[11pt]{ltjsarticle}
\usepackage{amsmath,amssymb,mathtools}
\usepackage{physics}
\usepackage{tikz-cd}
\usepackage{hyperref}
\usepackage{shortvrb} 
\hypersetup{colorlinks=true,linkcolor=blue,urlcolor=teal}

\title{構造塔と最小層の学部生向けガイド}
\author{Codex (OpenAI) --- TeXファイル生成\\Codex (OpenAI) --- Leanファイル生成}
\date{2025年11月7日}

\begin{document}
\maketitle

\section*{はじめに}
本ノートは、\texttt{HierarchicalStructureTower.lean} に実装された \emph{StructureTowerWithMin} の具体例を、学部生でも追える形で解説する。構造塔 (structure tower) とは、ある集合 $X$ の上に層 (layer) と呼ばれる部分集合族 $\{X_i\}_{i\in I}$ を載せ、各点 $x\in X$ に「最小層」$\operatorname{minLayer}(x)$ を割り当てるデータである。以下では、離散位相・フィルター・主イデアル・崩壊構造という4種類の例を順に扱い、最後に簡単な図式を示す。

\section{離散位相による自明な最小層}
\subsection*{設定}
\begin{itemize}
  \item 基礎集合: 任意の型 $X$。
  \item 層: $I = \mathcal{P}(X)$ とし、各 $U \subseteq X$ をそのまま層とみなす。
  \item 順序: 包含順序 $U \le V \iff U \subseteq V$。
  \item 最小層: 各 $x \in X$ に単集合 $\{x\}$ を割り当てる。
\end{itemize}
離散位相空間では全ての単集合が開であるため、$\{x\}$ は $x$ を含む最小の開集合であり、minLayer の候補として自然である。

\subsection*{数学的ポイント}
\begin{enumerate}
  \item \textbf{被覆性}: 任意の $x$ に対し $x \in X = \text{layer}(X)$ が成り立つ。
  \item \textbf{単調性}: 含意 $U \subseteq V \Rightarrow U \le V$ は明らか。
  \item \textbf{最小性}: $x \in U$ なら $\{x\} \subseteq U$、よって $\operatorname{minLayer}(x) \le U$。
\end{enumerate}

\section{フィルター塔と収束の視点}
\subsection*{設定}
\begin{itemize}
  \item 基礎集合: 位相空間 $(X,\mathcal{T})$。
  \item 層: フィルターの集合 $\mathsf{Filt}(X)$ の順序双対 $\mathsf{Filt}(X)^{\mathrm{op}}$。
  \item 各層: $F \in \mathsf{Filt}(X)^{\mathrm{op}}$ に対し $\text{layer}(F) = \{x \mid F \le \mathcal{N}(x)\}$。
  \item 最小層: $\operatorname{minLayer}(x) = \mathcal{N}(x)$ (近傍フィルター) の双対元。
\end{itemize}
フィルターは「情報の細かさ」を測る道具であり、より細かいフィルターほど塔の上位に位置する。近傍フィルターは $x$ での収束挙動を端的に記述するため、最小層として canonically 選べる。

\subsection*{メモ}
双対を取るのは「上に行く=細かくなる」という直感と順序を一致させるためである。$F \le G$ (元の順序) なら $G$ のほうが細かいが、双対を通すことで $F \le_{\mathrm{op}} G$ は逆向きとなり、塔の条件 (上方向へ単調) を満たしやすくなる。

\section{主イデアル塔}
\subsection*{設定}
\begin{itemize}
  \item 基礎集合: 可換環 $R$。
  \item 層: イデアル全体 $\operatorname{Ideal}(R)$ を包含順序付き集合として利用。
  \item 最小層: 各元 $r$ に対し主イデアル $(r) = \langle r \rangle = \operatorname{Ideal.span}\{r\}$。
\end{itemize}
この塔では代数的生成の概念が直接 minLayer に反映される。$r$ を含む任意のイデアル $I$ に対して $(r) \subseteq I$ なので、最小性が保証される。

\section{崩壊構造塔}
\subsection*{設定}
\begin{itemize}
  \item 基礎集合: 任意の集合 $X$。
  \item 添字集合: 任意の前順序付き集合 $(I,\le)$。
  \item 固定要素: $i_0 \in I$。
  \item 層: $\text{layer}(i) = \begin{cases}X & \text{if } i_0 \le i, \\ \varnothing & \text{otherwise}\end{cases}$。
  \item 最小層: 全て $i_0$ に固定。
\end{itemize}
全ての点が同じ最小層を共有する極端例であり、塔の構造が完全に「崩壊」している様子を示す。これは minLayer の自由度を議論する際のベースケースとして役立つ。

\section{図式による整理}
以下の図式では、離散位相塔と主イデアル塔の最小層選択が、構造塔圏における射としてどのようにつながるかを示す。
\[
\begin{tikzcd}[column sep=large]
  (X, \mathcal{P}(X)) \ar[r, "\operatorname{id}"] \ar[d, "f"'] & (X, \mathcal{P}(X)) \ar[d, "g"] \\
  (X, \mathsf{Filt}(X)^{\mathrm{op}}) \ar[r, "h"] & (R, \operatorname{Ideal}(R))
\end{tikzcd}
\]
ここで $f$ は各点をその近傍フィルターに送る射、$g$ は恒等射、$h$ は (例えば多項式写像を通じて) 点の像を生成する主イデアルを対応させる射を表す。矢印が可換であることは、minLayer の自然性 (Lean ファイル内の \texttt{minLayer\_preserving}) を意味する。

\section*{まとめと今後の学習}
\begin{itemize}
  \item 離散位相塔は「自明な例」、フィルター塔は「極端に密な例」、主イデアル塔は「代数的な正準例」、崩壊塔は「病的な極限例」である。
  \item Lean での形式化を通じて、最小層の存在と一意性が圏論的射の自然性を保証することが確認できる。
  \item 発展として、測度論的塔 (σ-代数)、複雑度クラス塔などを同様の枠組みで調べるとよい。
\end{itemize}

\end{document}
